Пусть $\xi_1, \xi_2...$ - НОР $P(\xi_1 = 1) = P(\xi_1 = -1) = \frac{1}{2}$
\begin{definition}
$S_0 = 0, ..., S_n = \xi_1 + .. + \xi_n$ - простейшее симметричное случайное блуждание
\end{definition}
\begin{definition} Зафиксируем $\omega \in \Omega$, то
$\{S_n(\omega), n \in \Z_{+}\}$ - траектория
\end{definition}
$P(S_n = k) = (\frac{1}{2})^n C_n^x$,\;\; где $x - (n-x) = k \;\Rightarrow\; x = \frac{n+k}{2} \;\Rightarrow\; P(S_n = k) =(\frac{1}{2})^n C_n^{\frac{n+k}{2}}$
\\
\\
\textbf{Задача:}  Пусть k > 0. $P(S_1 > 0, S_2 > 0,..., S_n = k) = ?$
Для начала заметим, что первым шагом мы в любом случае должны пойти наверх, то есть начальная точка 1 а не 0. Найдем обратную величину: кол-во траекторий, пересекающих ноль. Между каждой такой траекторией и траекторией из 1 в $-k$ можно задать биективное отображение (просто отражаем относительно первой точки пересечения с нулем).\\ Посчитаем количество траекторий, которые ведут в $(n, -k), S_1 > 0$ = количество тракторий в  $(n, k)$, пересекающих 0, $ S_1 > 0$ =  количество траекторий из (1, 1) в  $(n, -k)$ =  $C_{n-1}^{\frac{n-k - 2}{2}}$, откуда  $P(S_1 > 0, S_2 > 0,..., S_n = k) = (\frac{1}{2})^n (C_{n-1}^{\frac{n+k-2}{2}} - C_{n-1}^{\frac{n-k-2}{2}}) = (\frac{1}{2})^n (C_{n-1}^{\frac{n+k}{2}} - C_{n-1}^{\frac{n-k}{2}})$